%--- iliad ---
% title: Lecture 1: latent models and their conditions
%--- end ---
% Lecture 1. Uses the clustering illustration in fig/.
\documentclass[aspectratio=169,12pt]{beamer}
\usepackage[T1]{fontenc}
\usepackage{lmodern}
\usepackage{amsmath,amssymb}
\usetheme{default}
\usefonttheme{professionalfonts}
\setbeamercolor{normal text}{fg=black,bg=white}
\setbeamercolor{structure}{fg=black}
\setbeamercolor{frametitle}{fg=black,bg=white}
\setbeamerfont{frametitle}{size=\Large,series=\bfseries}
\setbeamertemplate{navigation symbols}{}
% Default beamer sets the frame title flush to the top-left, outside the text
% block: it started 3mm from the left edge while body text starts at 9mm. Align
% it with the body and give it room to breathe above.
\setbeamertemplate{frametitle}{%
  \vskip5mm\hspace*{2mm}%
  \usebeamerfont{frametitle}\usebeamercolor[fg]{frametitle}\insertframetitle%
  \par\vskip3mm}
\setbeamertemplate{footline}[frame number]
\setbeamertemplate{itemize items}[circle]
\setbeamersize{text margin left=9mm,text margin right=9mm}
\newcommand{\PI}{\mathcal P^+(I)}
\newcommand{\Up}{\mathord{\supseteq}}
\title{Condensation}
\subtitle{Models and conditions}
\author{Satya Benson}
\date{ILIAD, 21 September 2026}
\begin{document}

\begin{frame}{Condensation}
In empirical work, people look for features in AI systems which correspond
robustly to human concepts. In general, these can be hard to find.
\medskip

One ambitious approach to this problem is to think through from first
principles how agents in general might organize information about the world
into concepts.
\medskip

How should information be broken into reusable pieces? Is there one way to do
this which many different kinds of agents will converge on?
\medskip

We will propose two conditions that we think are normative for agents in
general: conditions an agent will want its own organization of information to
satisfy. Under those conditions, corresponding latent collections in two models
determine one another.
\end{frame}

\begin{frame}{Three overlapping observations}
Let $S,P,Q$ be independent fair bits.
\[
X_1=(S,P),\qquad X_2=(S,Q),\qquad X_3=(P,Q).
\]
One latent model assigns
\[
Y_{\{1,2\}}=S,\qquad Y_{\{1,3\}}=P,\qquad Y_{\{2,3\}}=Q.
\]
Each observation uses two latents. Each latent is recoverable from either
observation to which it is assigned.
\end{frame}

\begin{frame}{A latent model}
Fix finite-valued observables $(X_i)_{i\in I}$, with $I=\{1,\ldots,n\}$, and
write $\PI$ for the nonempty subsets of $I$.
\medskip

Each $A\in\PI$ is an \emph{address}, carrying a latent $Y_A$ assigned to the
observations in $A$. For a family $\mathcal F\subseteq\PI$ of addresses, write
$Y_{\mathcal F}=(Y_A)_{A\in\mathcal F}$ for the corresponding tuple of latents.
\medskip

The joint law of $(X,Y)$ must satisfy the \textbf{latent variable model
condition (LVM)}: each observation is determined by the latents assigned to it,
\[
H\bigl(X_i\mid Y_{\{A\,:\,i\in A\}}\bigr)=0\qquad(i\in I).
\]
Unused slots are constant. This condition does not assert that an agent knows
all the realized latent values.
\end{frame}

% Exactly two information-theory refresher frames follow.
\begin{frame}{Refresher: entropy and recovery}
All variables are finite-valued. Logarithms have base two.
\[
H(U)=-\sum_u p(u)\log p(u),\qquad
H(U\mid V)=\sum_v p(v)H(U\mid V=v).
\]
\[
H(U\mid V)=0
\quad\Longleftrightarrow\quad
U=f(V)\text{ almost surely, for some }f.
\]
The chain rule and conditioning inequality:
\[
\begin{aligned}
H(U,V\mid W)&=H(U\mid W)+H(V\mid U,W),\\
H(U\mid V,W)&\le H(U\mid W).
\end{aligned}
\]
\end{frame}

\begin{frame}{Refresher: dependence}
\[
I(U;V\mid W)=H(U\mid W)-H(U\mid V,W)\ge0.
\]
This is zero exactly when $U$ and $V$ are conditionally independent given $W$.
Omitting $W$ gives $I(U;V)$.
\medskip

For laws $p,q$ on the same finite set,
\[
D_{\mathrm{KL}}(p\Vert q)=\sum_u p(u)\log\frac{p(u)}{q(u)}\ge0.
\]
KL divergence appears in the score interpretation. The correspondence proof
uses the chain rule and nonnegativity.
\end{frame}

\begin{frame}{LVM and reconstruction}
\textbf{LVM:} the assigned latents determine the answer.
\[
H\bigl(X_i\mid Y_{\{A\,:\,i\in A\}}\bigr)=0.
\]
\textbf{Reconstruction:} each contributing observation recovers the slot.
\[
H(Y_A\mid X_i)=0\qquad\text{whenever }i\in A.
\]
An address alone does not imply the second condition.
\medskip

In our example, $Y_{\{1,2\}}=S$ is recoverable from both $(S,P)$ and $(S,Q)$.
\end{frame}

\begin{frame}{Upward-closed families}
A family $\mathcal F\subseteq\PI$ is \textbf{upward closed} if
\[
A\in\mathcal F,\quad A\subseteq B
\quad\Longrightarrow\quad B\in\mathcal F.
\]
Write $\Up A=\{C\in\PI:A\subseteq C\}$ for the addresses containing $A$; every
such family is upward closed. For $I=\{1,2,3\}$,
\[
\begin{aligned}
\Up\{1\}&=\{\{1\},\{1,2\},\{1,3\},\{1,2,3\}\},\\
\Up\{2\}&=\{\{2\},\{1,2\},\{2,3\},\{1,2,3\}\}.
\end{aligned}
\]
Their intersection is $\{\{1,2\},\{1,2,3\}\}$.
\medskip

These are families of addresses. $Y_{\mathcal F}$ is a tuple of variables.
\end{frame}

\begin{frame}{The Markov condition}
For every pair of upward-closed families $\mathcal F,\mathcal G\subseteq\PI$, require
\[
I\bigl(Y_{\mathcal F};Y_{\mathcal G}
\mid Y_{\mathcal F\cap\mathcal G}\bigr)=0.
\]
Given the latents at the addresses in the intersection $\mathcal F\cap\mathcal G$,
the two collections are conditionally independent.
\medskip

For approximate models, record the largest defect:
\[
\eta_Y=\max_{\mathcal F,\mathcal G\ \mathrm{upward\ closed}}
I\bigl(Y_{\mathcal F};Y_{\mathcal G}
\mid Y_{\mathcal F\cap\mathcal G}\bigr).
\]
\end{frame}

\begin{frame}{A hierarchical model of 36 labelled points}
\centering
\includegraphics[width=0.98\linewidth,height=5.0cm,keepaspectratio]{fig/hierarchical-points.pdf}

\end{frame}

\begin{frame}{Why we ask for the conditions only approximately}
The cluster center specifies a distribution for its points, not the points
themselves, so a point does not determine the center it was drawn from:
\[
H(M_{gc}\mid X_i)>0.
\]
Every kernel here has full support: no finite set of points pins a parameter.
\medskip

LVM and the Markov condition hold exactly here. The reconstruction condition
holds only approximately.
\medskip

So we ask for reconstruction up to a defect: for $i\in A$,
\[
H(Y_{\Up A}\mid X_i)\le\delta ,
\]
and prove a correspondence whose error grows with $\delta$ and $\eta$.
\end{frame}

\begin{frame}{Two ways the conditions can fail}
Return to $X_1=(S,P), X_2=(S,Q), X_3=(P,Q)$ with independent fair bits.
\medskip

\textbf{Copy each observation into its singleton slot.}
LVM and reconstruction both hold, but the common slots are constant:
\[
I(Y_{\Up\{1\}};Y_{\Up\{2\}}\mid Y_{\Up\{1,2\}})=1.
\]
\textbf{Put $(S,P,Q)$ only in the top slot.}
The Markov condition holds, but
\[
H(Y_{\{1,2,3\}}\mid X_1)=H(S,P,Q\mid S,P)=1.
\]
The first model hides shared information below constant common slots.
The second assigns information to a slot that one observation cannot recover.
\end{frame}



\begin{frame}{Perfect condensation}
A latent model is a perfect condensation precisely when
\begin{itemize}
\item the Markov condition holds exactly,
\item reconstruction holds: $H(Y_A\mid X_i)=0$ whenever $i\in A$.
\end{itemize}
\medskip
LVM is already part of being a latent model.
\medskip

The combination constrains where information can be placed.
Not every observable law admits a perfect condensation.
\end{frame}

\begin{frame}{The collections we will compare}
The tuple $Y_{\Up A}$ contains the slot at $A$ and all its supersets.
For example,
\[
Y_{\Up\{1,2\}}=(Y_{\{1,2\}},Y_{\{1,2,3\}}).
\]
Our theorem compares $Y_{\Up A}$ and $Z_{\Up A}$ in two models.
Individual slots may contain different amounts of duplicated information.
\medskip

Note that if we are using the latent at $A$ to predict the observations in
$A$, every latent at a $C\supseteq A$ is assigned to those observations too,
so it is already part of what we would need. A family we would actually use to predict an observable is
closed upward.
\end{frame}

\end{document}
